\SourcePage{456}{459}
\section{Irreducible representations of point groups}

We now turn to the explicit determination of the irreducible representations of point groups. The overwhelming majority of molecules possess only symmetry axes of orders two, three, four, and six. We shall therefore omit the icosahedral groups $Y$, $Y_h$; consider $C_n$, $C_{nh}$, $C_{nv}$, $D_n$, $D_{nh}$ only for $n=1,2,3,4,6$; and consider $S_{2n}$, $D_{nd}$ only for $n=1,2,3$.

The characters of the representations of these groups are given in Table~7. Isomorphic groups have the same representations and are placed together in one array. Numbers preceding the group-element symbols in the first rows give the numbers of elements in the corresponding classes (see \S~93). The first columns contain the conventional labels for the representations. One-dimensional representations are denoted by $A$, $B$, two-dimensional ones by $E$, and three-dimensional ones by $F$ (the symbol $E$ for a two-dimensional irreducible representation must not be confused with $E$ for the identity element!)\footnote{The reason why two complex-conjugate one-dimensional representations are denoted as a single two-dimensional representation will become clear in \S~96.}. Basis functions of the $A$ representations are symmetric, and those of the $B$ representations antisymmetric, under rotations about the principal axis of order $n$. Functions having different symmetry under reflection $\sigma_h$ are distinguished by the number of primes (one or two), while the subscripts $g$ and $u$ specify symmetry under inversion. Alongside the representation labels, the letters $x,y,z$ indicate the representation according to which the coordinates themselves transform; the $z$ axis is everywhere chosen along the principal symmetry axis. The letters $\varepsilon$ and $\omega$ mean
\[
\varepsilon=e^{2\pi i/3},\qquad \omega=e^{2\pi i/6}=-\omega^4,
\qquad \varepsilon+\varepsilon^2=-1,\qquad \omega^2-\omega=-1.
\]

\begin{table}[p]
\centering\small
\caption{Characters of irreducible representations of point groups}
\begin{tabular}{c|cc|c|cc|c|ccc}
$C_i$&$E$&$I$&$C_2$&$E$&$C_2$&$C_s$&$E$&$\sigma$\\\hline
$A_g$&1&1&$A;z$&1&1&$A';x,y$&1&1\\
$A_u;x,y,z$&1&-1&$B;x,y$&1&-1&$A'';z$&1&-1
\end{tabular}

\medskip
\begin{tabular}{c|rrrr|c|rrrr|c|rrrr}
$C_{2h}$&$E$&$C_2$&$\sigma_h$&$I$&$C_{2v}$&$E$&$C_2$&$\sigma_v$&$\sigma'_v$&$D_2$&$E$&$C_2$&$C'_2$&$C''_2$\\\hline
$A_g$&1&1&1&1&$A_1;z$&1&1&1&1&$A$&1&1&1&1\\
$B_g$&1&-1&-1&1&$B_2;y$&1&-1&-1&1&$B_3;x$&1&-1&-1&1\\
$A_u;z$&1&1&-1&-1&$A_2$&1&1&-1&-1&$B_1;z$&1&1&-1&-1\\
$B_u;x,y$&1&-1&1&-1&$B_1;x$&1&-1&1&-1&$B_2;y$&1&-1&1&-1
\end{tabular}

\medskip
\begin{tabular}{c|rrr|c|rrr}
$C_{3v}$&$E$&$2C_3$&$3\sigma_v$&$D_3$&$E$&$2C_3$&$3U_2$\\\hline
$A_1;z$&1&1&1&$A_1$&1&1&1\\
$A_2$&1&1&-1&$A_2;z$&1&1&-1\\
$E;x,y$&2&-1&0&$E;x,y$&2&-1&0
\end{tabular}
\end{table}

\SourcePage{457}{460}
\begin{table}[p]\centering\small
\caption*{Table 7 (continued)}
\begin{tabular}{c|rrrrrr}
$C_6$&$E$&$C_6$&$C_3$&$C_2$&$C_6^5$&$C_6^4$\\\hline
$A;z$&1&1&1&1&1&1\\ $B$&1&-1&1&-1&1&-1\\
$E_1$&1&$\omega$&$\omega^2$&-1&$-\omega$&$-\omega^2$\\
&1&$-\omega$&$\omega^2$&-1&$\omega$&$-\omega^2$\\
$E_2;x\pm iy$&1&$\omega^2$&$-\omega$&1&$\omega^2$&$-\omega$\\
&1&$-\omega^2$&$-\omega$&1&$-\omega^2$&$\omega$
\end{tabular}

\medskip
\begin{tabular}{c|rrrrr}
$C_{4v},D_4,D_{2d}$&$E$&$C_2$&$2C_4$&$2U_2$&$2U'_2$\\\hline
$A_1$&1&1&1&1&1\\ $A_2$&1&1&1&-1&-1\\
$B_1$&1&1&-1&1&-1\\ $B_2$&1&1&-1&-1&1\\ $E;x,y$&2&-2&0&0&0
\end{tabular}
\end{table}

\SourcePage{458}{461}
\begin{table}[p]\centering\small
\caption*{Table 7 (continued)}
\begin{tabular}{c|rrrrrr}
$D_6,C_{6v},D_{3h}$&$E$&$C_2$&$2C_3$&$2C_6$&$3U_2$&$3U'_2$\\\hline
$A_1$&1&1&1&1&1&1\\ $A_2$&1&1&1&1&-1&-1\\
$B_1$&1&-1&1&-1&1&-1\\ $B_2$&1&-1&1&-1&-1&1\\
$E_2$&2&2&-1&-1&0&0\\ $E_1;x,y$&2&-2&-1&1&0&0
\end{tabular}

\medskip
\begin{tabular}{c|rrrr|c|rrrrr}
$T$&$E$&$3C_2$&$4C_3$&$4C_3^2$&$O,T_d$&$E$&$8C_3$&$3C_2$&$6\sigma_d$&$6S_4$\\\hline
$A$&1&1&1&1&$A_1$&1&1&1&1&1\\
$E$&1&1&$\varepsilon$&$\varepsilon^2$&$A_2$&1&1&1&-1&-1\\
&1&1&$\varepsilon^2$&$\varepsilon$&$E$&2&-1&2&0&0\\
$F;x,y,z$&3&-1&0&0&$F_2;x,y,z$&3&0&-1&1&-1\\
&&&&&$F_1$&3&0&-1&-1&1
\end{tabular}
\end{table}

Clearly, when $\widehat G$ acts on a basis function $\psi$, the latter can only be multiplied by a $g$th root of unity, i.e.\footnote{For the point group $C_n$, possible choices are, for example, $\psi=e^{ik\varphi}$ ($k=1,2,\ldots,n$), where $\varphi$ is the angle of rotation about the axis measured from some fixed direction.}
\[
\widehat G\psi=e^{2\pi ik/g}\psi,\qquad k=1,2,\ldots,g.
\]
The group $C_{2h}$ (and the groups $C_{2v}$ and $D_2$ isomorphic to it) is Abelian, so all its irreducible representations are one-dimensional; their characters can only be $\pm1$ (since the square of every element is $E$).

Consider next the group $C_{3v}$. Relative to $C_3$, it also contains reflections $\sigma_v$ in vertical planes (all belonging to one class). A function invariant under rotation about the axis (a basis function of the $A$ representation of $C_3$) may be symmetric or antisymmetric under the reflections $\sigma_v$. Functions multiplied by $\varepsilon$ and $\varepsilon^2$ under a $C_3$ rotation (basis functions of the complex-conjugate $E$ representations) are interchanged by reflection\footnote{They may be chosen, for example, as $\psi_1=e^{i\varphi}$, $\psi_2=e^{-i\varphi}$. Reflection in a vertical plane changes the sign of $\varphi$.}. It follows that $C_{3v}$ (and the isomorphic group $D_3$) has two one-dimensional irreducible representations and one two-dimensional irreducible representation with the characters displayed in the

\SourcePage{459}{462}
table. That these are all the irreducible representations follows from $1^2+1^2+2^2=6$, equal to the order of the group.

The characters of representations of the other groups of the same type ($C_{4v}$, $C_{6v}$) are obtained by analogous arguments.

The group $T$ is obtained from $D_2\equiv V$ by adjoining rotations about four inclined axes of order three. A function invariant under transformations of $V$ (the basis of representation $A$) may be multiplied by $1$, $\varepsilon$, or $\varepsilon^2$ under a $C_3$ rotation. The basis functions of the three one-dimensional representations $B_1$, $B_2$, $B_3$ of $V$ are interchanged by rotations about the third-order axes (as is evident if the coordinates $x,y,z$ themselves are chosen as these functions). We thus obtain three one-dimensional and one three-dimensional irreducible representations ($1^2+1^2+1^2+3^2=12$).

Finally, consider the isomorphic groups $O$ and $T_d$. The group $T_d$ is obtained from $T$ by adjoining reflections $\sigma_d$ in planes, each of which passes through two third-order axes. A basis function of the identity representation $A$ of $T$ may be symmetric or antisymmetric under these reflections (all of which form one class), giving two one-dimensional representations of $T_d$. Functions multiplied by $\varepsilon$ or $\varepsilon^2$ under rotation about a third-order axis (the bases of the complex-conjugate $E$ representations of $T$) are interchanged under reflection in a plane containing that axis, so that a single two-dimensional representation results.

Of the three basis functions of representation $F$ of $T$, one transforms into itself under reflection (and may either remain unchanged or change sign), while the other two are interchanged. Altogether this gives two one-dimensional, one two-dimensional, and two three-dimensional representations\footnote{Irreducible representations of higher dimensions (four and five) occur in the icosahedral groups.}.

The representations of the remaining point groups of interest follow directly from those already written down, since these groups are direct products of groups already considered with $C_i$ (or $C_s$). Specifically,
\[
\begin{gathered}
C_{3h}=C_3\times C_s,\quad D_{2h}=D_2\times C_i,\quad D_{3d}=D_3\times C_i,\quad O_h=O\times C_i,\\
C_{4h}=C_4\times C_i,\quad D_{4h}=D_4\times C_i,\quad D_{6h}=D_6\times C_i,\\
C_{6h}=C_6\times C_i,\quad S_6=C_3\times C_i,\quad T_h=T\times C_i.
\end{gathered}
\]

\SourcePage{460}{463}
Each of these direct products has twice as many irreducible representations as the original group; half are symmetric (denoted by the subscript $g$) and half antisymmetric (subscript $u$) under inversion.

Their characters are obtained from the characters of the original group's representations by multiplication by $\pm1$ (in accordance with (94.24)). Thus, for $D_{3d}$ we obtain
\[
\begin{array}{c|rrrrrr}
D_{3d}&E&2C_3&3U_2&I&2S_6&3\sigma_d\\\hline
A_{1g}&1&1&1&1&1&1\\ A_{2g}&1&1&-1&1&1&-1\\ E_g&2&-1&0&2&-1&0\\
A_{1u}&1&1&1&-1&-1&-1\\ A_{2u}&1&1&-1&-1&-1&1\\ E_u&2&-1&0&-2&1&0
\end{array}
\]
